\section*{Three languages for one bound: a reading guide}
\addcontentsline{toc}{section}{Three languages for one bound: a reading guide}

\noindent\motiv{\vet{Lieb-Robinson bounds are written in two dialects. The review~\cite{chen-lucas-yin}
speaks of a graph whose edges carry Hamiltonian terms $H_e$, of a \emph{site matrix} $h_{uv}$
adding up the strengths of the terms containing two sites, and of walks on the graph; the
mathematical-physics papers~\cite{nachtergaele-sims,nsy2019} speak of an interaction $\Phi$
assigning a self-adjoint term $\Phi(Z)$ to each finite set $Z$ of sites, of a \emph{decay
function} $F$ with a reproducing property, and of the $F$-norm $\norm{\Phi}_F$
(Chapter~\ref{sec:closed} defines the site matrix and the decay functions).

These notes use a third, slightly more general language, chosen so that one theorem covers Hamiltonian and open
dynamics alike: a finite family of \emph{local generators} $L_i$ with supports $S_i$ and
strengths $\kappa_i$, and \emph{paths} of supports from the probe's region to the observable's.
The table says which object plays which role, and where each conversion is proved.}}

\medskip
\noindent\vet{\begin{longtable}{@{}>{\raggedright\arraybackslash}p{0.29\linewidth}>{\raggedright\arraybackslash}p{0.30\linewidth}>{\raggedright\arraybackslash}p{0.35\linewidth}@{}}
\toprule
Review (graph language) & Mathematical physics (interactions) & These notes (generators) \\
\midrule
\endfirsthead
\toprule
Review & Mathematical physics & These notes \\
\midrule
\endhead
edge Hamiltonian $H_e$, $e=\{u,v\}$ & interaction term $\Phi(Z)$, $Z$ a finite set &
generator $L_i$ with support $S_i$; Hamiltonian case $L_i=\ii\comm{\Phi(S_i)}{\cdot}$
(Chapter~\ref{sec:qudit}) \\
strength $\norm{H_e}$ & $\norm{\Phi(Z)}$ & constant $\kappa_i$; equals $\norm{\Phi(S_i)}$ in the
Hamiltonian case; \vetnew{for a Lindblad term the direct strength
$\norm H+\sum_j\norm{\Gamma_j}^2$, by a commutator identity, or $\norm{L_S}$ through the
conditional expectation} (Chapters~\ref{sec:qudit}, \ref{sec:approx}) \\
Heisenberg evolution $e^{\ii tH}Ae^{-\ii tH}$ & $\tau_t(A)$, the automorphism group &
$e^{tL}A$ with $L=\sum_i L_i$; hypotheses G1--G3 (Definition~\ref{def:local-generators}); equal
to the conjugation in the Hamiltonian case (Chapter~\ref{sec:qudit}) \\
operators on the qubits of a region & \vetnew{$\cA_X=\bigotimes_{x\in X}\Mat_n\otimes1$}, the tensor
factor & the matrices local to $X$: zero between configurations differing off $X$, entries
depending on the letters inside $X$ (Chapter~\ref{sec:qudit}) \\
\vetnew{bosons absent} & \vetnew{Weyl operators $W(f)$ with $\tau_t(W(f))=W(f_t)$, the kernel $f_t$ by
normal modes} & \vetnew{a quadratic generator $K$ on the Weyl labels, the flow $T_t=e^{tK}$ and
$\alpha_t(W(f))=W(T_tf)$; the site matrix is $|K|$ (Chapter~\ref{sec:harmonic})} \\
$D(X,Y)$ absent & \vetnew{$D(X,Y)=\min\{\sum_{X\times\bd Y},\sum_{\bd X\times Y}\}$}, eq.~(3.24) of
\cite{nsy2019} & the smaller of the two \vetnew{walk sums or sums of the decay function}, by the role exchange of
automorphisms (Chapter~\ref{sec:qudit}) \\
site matrix $h_{uv}=\sum_{e\ni u,v}\norm{H_e}$, eq.~(3.35) & $\sum_{Z\ni u,v}\norm{\Phi(Z)}\le
\norm{\Phi}_F\,F(\dist(u,v))$, eq.~(3.14) & site matrix $h_{uv}=\sum_{i:\,u,v\in S_i}\kappa_i$, with
diagonal $h_{uu}=\sum_{i:\,u\in S_i}\kappa_i$ (Chapter~\ref{sec:closed}) \\
walk on the graph; irreducible path (Chen--Lucas) & chain of interaction supports with
boundary-crossing sets, eqs.~(3.41)--(3.42) of~\cite{nsy2019} & path: indices $i_1,\dots,i_k$
with $S_{i_1}$ crossing $Y$ and each later support crossing the previous one (no term follows
itself, none is nested in its predecessor), with first contact with $X$ at step $k$
(Definition~\ref{def:path}) \\
\vetnew{$\mathcal C_{X,Y}(t)\le\sum_{x\in X,y\in Y}(e^{2|t|h})_{xy}$}, eq.~(3.36), with
$\mathcal C=\tfrac12\sup\norm{\comm{A(t)}B}$ by its eq.~(3.24) (its (3.22) omits the $\tfrac12$, a
misprint) &
$\norm{\comm{\tau_t(A)}{B}}\le\frac{2\norm A\norm B}{C_F}(e^{2\norm{\Phi}_F C_F t}-1)D(X,Y)$,
Thm.~3.1 of~\cite{nsy2019} & path series (Theorem~\ref{thm:master}), then the two closed forms
as corollaries (Chapter~\ref{sec:closed}); \vetnew{the matrix form is the review's (3.36) with
its constant doubled and the $-1$ that the series proof gives} \\
Haar average over unitaries on $X^c$, Prop.~4.1 & conditional expectation $\mathbb E_X$ as a
finite group average, eq.~(4.6) of~\cite{nsy2018} & conditional expectation $P_X$, the
normalized partial trace, computed as a finite twirl over the Weyl unitaries of $X^c$
(Chapter~\ref{sec:approx}); \vetnew{also a second route to G1 for Lindblad terms, with the
strength $\norm{L_S}$} \\
$\norm{\mathbb E_SA-A}\le\sup\norm{\comm A{B}}\le2\norm{\mathbb E_SA-A}$, Prop.~4.1 & Lemma~4.1 of~\cite{nsy2019} &
local approximation, the same constants (Theorem~\ref{thm:localapprox}); on any causal net the
converse holds for every element of $\cA(X)$ \\
Truncated dynamics $\norm{\tau_t(A)-\tau^{B}_t(A)}$ under the Hamiltonian restricted to $B$,
Prop.~4.3 & --- & $\norm{A(t)-P_{X'}A(t)}$ with the conditional expectation as the local
approximant, a corollary of the walk bound with our constants (Theorem~\ref{thm:truncated});
stated as $\varepsilon$-locality of $e^{tL}$, the approximate form of a cellular automaton's
locality (Corollary~\ref{cor:automaton}) \\
Lindblad generator, Thm.~3.14, with the norms of the terms as strengths & --- & the same
theorem: a Lindblad family is a generator family satisfying G1--G3, G3 as the contraction
hypothesis implied by the review's complete positivity (Chapter~\ref{sec:approx}) \\
--- & CAR algebra, parity automorphism, $\comm{A}{B}=0$ or $\{A,B\}=0$ for disjoint supports,
Prop.~2.1 of~\cite{nsy2018} & graded net; graded bracket $\gcomm{A}{B}$ (Chapter~\ref{sec:fermions}) \\
\bottomrule
\end{longtable}}


\medskip
\noindent\vet{\textbf{Reconnections.} \vetnew{Each chapter from Chapter~\ref{sec:master} on ends
with the sources' statements recovered in the sources' notation: the review's Theorem~3.4 in
Chapter~\ref{sec:master}; Theorem~3.1 of~\cite{nsy2019} in Chapters~\ref{sec:closed}
and~\ref{sec:qudit}; the review's bound for spin systems, its eq.~(3.36), in
Chapter~\ref{sec:qudit}; the review's Theorem~3.14 and Proposition~4.1 in
Chapter~\ref{sec:approx}; Theorem~3.1 of~\cite{nrss2009} in Chapter~\ref{sec:harmonic}; Theorem~3.1
of~\cite{nsy2018} in the chapter to come.} Chapter~\ref{sec:master} sketches the route of the
sources through the Jacobi identity.}
